\section{单周期、多周期与流水线的共同状态边界}

设一条指令要依次完成取指 IF、译码 ID、执行 EX、访存 MEM 和写回 WB。三种组织的差别不是指令语义，而是组合逻辑放在哪些时钟边界之间。

\begin{longtable}{L{0.2\linewidth}L{0.34\linewidth}L{0.36\linewidth}}
\toprule
组织 & 每拍发生什么 & 决定时钟周期的路径 \\
\midrule
单周期 & 一拍内走完 IF 到 WB & 最慢指令的整条组合路径 \\
多周期 & 一条指令分几拍复用 ALU 和存储端口 & 最慢阶段，加上阶段寄存器开销 \\
流水线 & 多条指令分别占据不同阶段 & 最慢流水阶段，加上流水寄存器开销 \\
\bottomrule
\end{longtable}

若五段组合延迟分别为 200、120、180、250、100 ps，流水寄存器的时钟到输出、建立时间和偏斜预算合计 30 ps，则单周期下限约为 $850$ ps，理想五级流水周期下限是 $250+30=280$ ps。流水化缩短的是连续指令的启动间隔，不会把单条指令的五段工作凭空删掉；单条指令还要跨过五个边界，延迟约为 $5\times280=1400$ ps。

\begin{keyidea}
流水线提高吞吐，不保证降低单条指令时延。每增加一级流水寄存器，都会付出寄存器延迟、时钟功耗、旁路复杂度和冲刷代价。
\end{keyidea}

\section{五级流水线的数据与控制随行}

IF/ID、ID/EX、EX/MEM、MEM/WB 四组流水寄存器不仅保存数据，还必须保存后续阶段需要的控制位。例如译码阶段产生的“是否写寄存器、写回来自 ALU 还是内存、是否读写数据存储器”，要与对应指令的数据一路同行。若控制位错位，机器可能用上一条指令的写使能提交下一条指令的结果。

\begin{lstlisting}
IF:  PC -> instruction memory -> instruction, PC+4
ID:  decode -> source values, immediate, destination id, controls
EX:  ALU/compare -> result, branch decision
MEM: load/store transaction -> read data or store completion
WB:  select result -> architectural register write
\end{lstlisting}

分支决定在 EX 才得到时，IF 和 ID 可能已经装入两条错误路径指令。控制器必须把它们变成不写任何状态的气泡，并把 PC 重定向到正确目标。等待内存时则不同：相关阶段和上游寄存器必须保持，不能一边等待旧事务一边让后续指令覆盖其地址、数据和控制位。

\begin{classicnote}
流水线图要同时画三类线：数据向前走，旁路结果逆着阶段边界送回 EX，停顿与冲刷控制横向冻结或清空流水寄存器。只画五个方框，无法证明机器正确。
\end{classicnote}
